%LTeX: language=fr-FR
\documentclass[12pt]{article}
\usepackage[margin=0.8in]{geometry}
\usepackage[utf8]{inputenc}
\usepackage{graphicx}
\usepackage[export]{adjustbox}
\usepackage{amsmath}
\usepackage{amsfonts}
\usepackage{amssymb}
\usepackage[version=4]{mhchem}
\usepackage{stmaryrd}
\usepackage{bbold}
%\usepackage[fallback]{xeCJK}
%\usepackage{polyglossia}
\usepackage{fontspec}
\usepackage{censor}
\usepackage{currfile}
%\setCJKmainfont{Noto Serif CJK TC}
\usepackage{enumitem}
\usepackage{tikz}
\usetikzlibrary {arrows.meta, quotes} 
\usepackage{tkz-tab}
\usepackage{pgfplots}
\pgfplotsset{compat=1.18}
\usepackage{amsthm}
\usepackage{tcolorbox}
\usepackage{array}
\usepackage{tabularray}
\usepackage{tabularx}
\usepackage{makecell}
\usepackage{esvect}
%\usepackage{changepage}   % for the adjustwidth environment
\usepackage{enumitem}
\usepackage{tasks}
\usepackage{fancyhdr}
\usepackage[french]{babel}
\usepackage{lastpage}

\tcbuselibrary{theorems}

\definecolor{GrayCen}{HTML}{d2d2d2}
\definecolor{GrayBG}{HTML}{d9d9d9}
\definecolor{BlueDef}{HTML}{104e8b}
\definecolor{GreenProp}{HTML}{025540}
\definecolor{RedMet}{HTML}{cd5556}
\definecolor{BlackSav}{HTML}{000000}
\def\censorcolor{GrayCen}
\let\svcensorrule\censorrule
\renewcommand\censorrule[1]{%
\textcolor{\censorcolor}{\svcensorrule{#1}}}

\newtcbtheorem
[]% init options
{definition}% name
{Définition}% title
{%
  colback=GrayBG,
  colframe=BlueDef,
  fonttitle=\bfseries,
}% options
{def}% prefix

\newtcbtheorem
[]% init options
{proposition}% name
{Proposition}% title
{%
  colback=GrayBG,
  colframe=GreenProp,
  fonttitle=\bfseries,
}% options
{prop}% prefix

\newtcbtheorem
[]% init options
{method}% name
{Méthode}% title
{%
  theorem name,
  colback=GrayBG,
  colframe=RedMet,
  fonttitle=\bfseries,
}% options
{met}% prefix

\newtcbtheorem
[]% init options
{savoir}% name
{Savoir-faire du chapitre}% title
{%
  theorem name,
  colback=GrayBG,
  colframe=BlackSav,
  fonttitle=\bfseries,
}% options
{sav}% prefix

\newtheoremstyle{break}
{\topsep}{\topsep}%
{}{}%
{\bfseries}{}%
{\newline}{}%
\theoremstyle{break}
\newtheorem*{example}{Exemple~:}
\newtheorem*{examples}{Exemples~:}
\newtheorem*{remark}{Remarque~:}
\newtheorem*{remarks}{Remarques~:}
% Variante sans retour à la ligne après « Exemple. »
\newtheoremstyle{inline}
{\topsep}{\topsep}%
{}{}%
{\bfseries}{}%
{ }{}%
\theoremstyle{inline}
\newtheorem*{example*}{Exemple~:}

\newenvironment{demo}[1][\proofname]{%
  \begin{proof}[#1]$ $\newline
  }{%
  \end{proof}
}

%\setmainfont{CMU Serif}name

\title{
  \vspace{4em}
  \centering
  \tcbset{colframe=BlackSav,colback=GrayBG}
  \tcbox[tikznode]{
    Chapitre 3 \\ Fonctions affines
  }
}

\author{}
\date{}

\pagestyle{empty}
\graphicspath{ {../Images/} }
\input{\currfiledir ../def.tex}
\def \Document {Cours}


\setlength\parindent{0pt}
%\setlist[itemize]{topsep=0pt}
\setlist{nolistsep}
%\setlist[itemize]{noitemsep}
%\setlist[enumerate]{noitemsep}

\tikzset{cross/.style={path picture={
  \draw[black]
    (path picture bounding box.south east)--(path picture bounding box.north west)
    (path picture bounding box.south west)--(path picture bounding box.north east);
}}}

\newcolumntype{C}[1]{>{\centering\arraybackslash}m{#1}}
\newcolumntype{L}[1]{>{\raggedright\arraybackslash}m{#1}}
\newcolumntype{N}{@{}m{0pt}@{}}

\newcolumntype{Y}{>{\centering\arraybackslash}X}

\fancypagestyle{firststyle}
{
  \fancyhf{}
  \fancyhead[L]{\Structure}
  \fancyhead[R]{\Niveau}
  \fancyfoot[L]{Alexandre Fernandez}
  \fancyfoot[R]{{\thepage\  / \pageref{LastPage}}}
   %\renewcommand{\headrulewidth}{0pt} % removes horizontal header line
}


\begin{document}

\maketitle
\thispagestyle{firststyle}

\tableofcontents

% \newcommand{\multihide}[1]{\mbox{\vphantom{#1}\colorbox{gray!50}{\parbox{\textwidth}{{#1}}}}}
%\newcommand{\multihide}[1]{\mbox{\vphantom{#1}\colorbox{gray!50}{\parbox{\textwidth}{{\phantom{#1}}}}}}
% On crée une boîte temporaire pour stocker le texte
\newsavebox{\hidebox}

% -----------
% Macros des trous : version élève si \ELEVE est défini à la compilation
% (recette « prof + élève »), sinon version prof.
\ifdefined\ELEVE
% Version élève : contenu invisible
\newcommand{\hideM}[1]{\mbox{\vphantom{$#1$}\smash{\colorbox{gray!50}{\phantom{$\displaystyle #1$}}}} }
\newcommand{\hideMW}[1]{\mbox{\vphantom{$#1$}\smash{\colorbox{white}{\phantom{$\displaystyle #1$}}}} }
\newcommand{\hide}[1]{\mbox{\vphantom{#1}\smash{\colorbox{gray!50}{\phantom{#1}}}}}
\NewEnviron{multihide}{%
  \colorbox{gray!50}{
    \begin{lrbox}{\hidebox}%
        \begin{minipage}{\linewidth}%
          \BODY
        \end{minipage}%
    \end{lrbox}%
    % On affiche un espace invisible ayant les dimensions exactes de la boîte
    \vrule width 0pt height \ht\hidebox depth \dp\hidebox\hskip\wd\hidebox%
  }
}
% Corrigé long dans un cadre blanc : zone vide de même hauteur
% Argument optionnel : espace ajouté sous le corrigé, ex. \begin{multihideW}[2em]
\NewEnviron{multihideW}[1][0pt]{%
  \par\noindent\fbox{%
    \begin{lrbox}{\hidebox}%
      \begin{minipage}{\dimexpr\linewidth-2\fboxsep-2\fboxrule}%
        \BODY\par\vspace{#1}%
      \end{minipage}%
    \end{lrbox}%
    \vrule width 0pt height \ht\hidebox depth \dp\hidebox\hskip\wd\hidebox%
  }
}
\newcommand{\hidetk}[1]{}
\else
% Version prof : contenu visible sur fond gris
\newcommand{\hideM}[1]{\mbox{\vphantom{$#1$}\smash{\colorbox{gray!50}{$\displaystyle #1$}}} }
\newcommand{\hideMW}[1]{\mbox{\vphantom{$#1$}\smash{\colorbox{white}{$\displaystyle #1$}}} }
\newcommand{\hide}[1]{\mbox{\vphantom{#1}\smash{\colorbox{gray!50}{#1}}}}
\NewEnviron{multihide}{%
  \colorbox{gray!50}{
    \begin{minipage}{\linewidth}%
      \BODY
    \end{minipage}%
  }
}
% Corrigé long dans un cadre blanc (argument optionnel : espace ajouté)
\NewEnviron{multihideW}[1][0pt]{%
  \par\noindent\fbox{%
    \begin{minipage}{\dimexpr\linewidth-2\fboxsep-2\fboxrule}%
      \BODY\par\vspace{#1}%
    \end{minipage}%
  }
}
\newcommand{\hidetk}[1]{#1}
\fi
% -----------

\newcommand{\wvec}[1]{\overrightarrow{#1}}
\newcommand{\ssi}[0]{\emph{ssi}\ }
\newcommand{\setR}[0]{$\mathbb{R}$}

\newpage

\pagestyle{fancy}

\fancyhf{} % clear existing header/footer entries
\fancyhead[R]{ {\Niveau}}
\fancyhead[L]{ {\Chapitre}}
%\fancyhead[R]{\large \textbf{\textbf{\Document\ (B) -- 1h}}}
\fancyfoot[L]{ A. Fernandez }
\fancyfoot[R]{{\thepage\  / \pageref{LastPage}}}

% Repère quadrillé : \repere{xmin}{xmax}{ymin}{ymax}
\newcommand{\repere}[4]{%
  \draw[very thin, gray!40] (#1,#3) grid (#2,#4);
  \draw[->] (#1,0) -- ({#2+0.4},0) node[below right] {\small $x$};
  \draw[->] (0,#3) -- (0,{#4+0.4}) node[above left] {\small $y$};
  \foreach \x in {#1,...,#2} {\ifnum\x=0\else\draw (\x,0.08) -- (\x,-0.08) node[below] {\scriptsize $\x$};\fi}
  \foreach \y in {#3,...,#4} {\ifnum\y=0\else\draw (0.08,\y) -- (-0.08,\y) node[left] {\scriptsize $\y$};\fi}
  \node[below left] at (0,0) {\scriptsize $0$};
}

\section{Fonctions affines}

\begin{minipage}[t]{0.55\linewidth}
  \begin{example}
  Une base de loisirs loue des vélos.
  Elle facture $5$~€ de frais d'inscription, puis $3$~€ par heure de location.
  On note $x$ le nombre d'heures de location et $f(x)$ le prix payé, en euros.

    \vspace{0pt}
    \begin{itemize}
      \item Pour $1$ heure, on paie : \hideM{3 \times 1 + 5 = 8} €.
      \item Pour $4$ heures, on paie : \hideM{3 \times 4 + 5 = 17} €.
      \item Pour $x$ heures, on paie : $f(x) = \hideM{3x + 5}$.
    \end{itemize}

    \medskip
    Le prix est la somme d'une partie qui \emph{dépend} du nombre d'heures ($3x$) et d'une partie \emph{fixe} ($5$).
    On dit que $f$ est une fonction affine.

    \medskip
    Sa représentation graphique est ci-contre.
  \end{example}
\end{minipage}
\hfill
\begin{minipage}[t]{0.4\linewidth}
  \vspace{0pt}
  \centering
  \begin{tikzpicture}[x=0.75cm, y=0.16cm]
    \draw[very thin, gray!40, xstep=1, ystep=5] (0,0) grid (7,30);
    \draw[->] (0,0) -- (7.6,0) node[right] {\small h};
    \draw[->] (0,0) -- (0,32) node[left] {\small prix (€)};
    \foreach \x in {1,...,7} \node[below] at (\x,0) {\scriptsize $\x$};
    \foreach \y in {5,10,...,30} \node[left] at (0,\y) {\scriptsize $\y$};
    \node[below left] at (0,0) {\scriptsize $0$};
    \draw[thick, domain=0:7] plot (\x,{3*\x+5});
    \draw[dashed] (1,0) -- (1,8) -- (0,8);
    \draw[dashed] (4,0) -- (4,17) -- (0,17);
    \fill (1,8) circle (2pt) (4,17) circle (2pt) (0,5) circle (2pt);
  \end{tikzpicture}
\end{minipage}

\begin{definition}{Fonction affine}{}
  \begin{multihide}
    Une fonction $f$ définie sur $\mathbb{R}$ est \textbf{affine} s'il existe deux nombres réels $a$ et $b$ tels que, pour tout réel $x$ :
    \[
      f(x) = ax + b.
    \]
    Cas particuliers :
    \begin{itemize}
      \item si $b = 0$, alors $f(x) = ax$ : on dit que $f$ est \textbf{linéaire}.% (situation de proportionnalité) ;
      \item si $a = 0$, alors $f(x) = b$ : on dit que $f$ est \textbf{constante}.
    \end{itemize}
  \end{multihide}
\end{definition}

\begin{examples}
  Compléter le tableau.

  \renewcommand{\arraystretch}{1.6}
  \begin{center}
  \begin{tabularx}{\linewidth}{|c|Y|Y|Y|}
    \hline
    Fonction & Affine ? & $a$ & $b$ \\
    \hline
    $f_1(x) = 2x - 3$ & \hide{oui} & \hideM{2} & \hideM{-3} \\
    \hline
    $f_2(x) = -x + 4$ & \hide{oui} & \hideM{-1} & \hideM{4} \\
    \hline
    $f_3(x) = 5x$ & \hide{oui, linéaire} & \hideM{5} & \hideM{0} \\
    \hline
    $f_4(x) = 7$ & \hide{oui, constante} & \hideM{0} & \hideM{7} \\
    \hline
    $f_5(x) = x^2 + 1$ & \hide{non} & & \\
    \hline
  \end{tabularx}
  \end{center}
\end{examples}

\begin{remark}
  Attention à l'ordre d'écriture : dans $f(x) = 4 - 0{,}5x$, on a $b = 4$ et le coefficient $a$ est le nombre qui multiplie $x$, c'est-à-dire $-0{,}5$.
\end{remark}

\newpage
\section{Représentation graphique}

\begin{proposition}{}{}
  \begin{multihide}
    Dans un repère, la représentation graphique d'une fonction affine $f(x) = ax + b$ est une \textbf{droite}.
    On dit que cette droite a pour \textbf{équation réduite} $y = ax + b$.
    \begin{itemize}
      \item $a$ est le \textbf{coefficient directeur} (ou \emph{pente}) de la droite ;
      \item $b$ est l'\textbf{ordonnée à l'origine} : la droite coupe l'axe des ordonnées au point $(0~;~b)$.
    \end{itemize}
  \end{multihide}
\end{proposition}

\begin{remarks}
  \begin{itemize}
    \item[]
    \item La droite d'une fonction linéaire ($b = 0$) passe par l'\hide{origine du repère}.
    \item La droite d'une fonction constante ($a = 0$) est \hide{horizontale}.
  \end{itemize}
\end{remarks}

\begin{method}{Tracer une droite}{}
  Pour tracer la droite d'équation $y = ax + b$, il suffit de connaître \textbf{deux points}.
  \begin{itemize}
    \item \textbf{Avec deux images :} on calcule les images de deux nombres, on place les deux points et on trace la droite.
    \item \textbf{Avec $a$ et $b$ :} on place le point $(0~;~b)$. À partir de ce point, on avance de $1$ vers la droite et on monte de $a$ (ou on descend si $a < 0$). On obtient un deuxième point.
  \end{itemize}
\end{method}

\begin{example}
  \begin{minipage}[t]{0.55\linewidth}
    \vspace{0pt}
    \textbf{1.} Tracer la droite $d_1$ de la fonction $f(x) = 2x - 1$.

    \renewcommand{\arraystretch}{1.4}
    \begin{center}
    \begin{tabular}{|c|C{1.5cm}|C{1.5cm}|}
      \hline
      $x$ & $0$ & $3$ \\
      \hline
      $f(x)$ & \hideM{-1} & \hideM{5} \\
      \hline
    \end{tabular}
    \end{center}

    Calcul : $f(3) = \hideM{2 \times 3 - 1 = 5}$.
  \end{minipage}
  \hfill
  \begin{minipage}[t]{0.5\linewidth}
    \vspace{0pt}
    \centering
    \begin{tikzpicture}[scale=0.6]
      \repere{-2}{5}{-2}{6}
      \hidetk{
        \draw[red, thick, domain=-0.5:3.5] plot (\x,{2*\x-1}) node[right] {$d_1$};
        \fill (0,-1) circle (4pt) (3,5) circle (4pt);
      }
    \end{tikzpicture}
  \end{minipage}

  \bigskip
  \begin{minipage}[t]{0.55\linewidth}
    \vspace{0pt}
    \textbf{2.} Tracer la droite $d_2$ de la fonction $g(x) = -0{,}5x + 3$.

    \medskip
    \begin{multihideW}[2em]
      On a $b = 3$ : on place le point $(0~;~3)$.

      \medskip
      On a $a = -0{,}5$ : on avance de $2$ et on descend de $1$.
    \end{multihideW}
  \end{minipage}
  \hfill
  \begin{minipage}[t]{0.5\linewidth}
    \vspace{0pt}
    \centering
    \begin{tikzpicture}[scale=0.6]
      \repere{-2}{5}{-2}{6}
      \hidetk{
        \draw[blue, thick, domain=-2:5] plot (\x,{-0.5*\x+3}) node[above] {$d_2$};
        \fill (0,3) circle (4pt) (2,2) circle (4pt);
        \draw[dashed, thick, ->] (0,3) -- (2,3);
        \draw[dashed, thick, ->] (2,3) -- (2,2);
      }
    \end{tikzpicture}
  \end{minipage}
\end{example}

\begin{method}{Lire graphiquement l'équation réduite d'une droite}{}
  \begin{enumerate}
    \item On lit $b$ : c'est l'ordonnée du point où la droite coupe l'axe des ordonnées.
    \item On lit $a$ : à partir d'un point de la droite, on avance de $1$ vers la droite et on compte de combien on monte ($a > 0$) ou on descend ($a < 0$) pour revenir sur la droite.
  \end{enumerate}
  Si la lecture n'est pas précise, on avance de plusieurs carreaux :
  \[
    a = \frac{\text{déplacement vertical}}{\text{déplacement horizontal}}.
  \]
\end{method}

\begin{example}
  \begin{minipage}[t]{0.45\linewidth}
    \vspace{0pt}
    \centering
    \begin{tikzpicture}[scale=0.55]
      \repere{-4}{6}{-4}{6}
      \draw[blue, thick, domain=-4:4] plot (\x,{\x+2}) node[above left] {$d_1$};
      \draw[red, thick, domain=-1:4] plot (\x,{-2*\x+4}) node[right] {$d_2$};
      \draw[green!50!black, thick, domain=-3:6] plot (\x,{2/3*\x-2}) node[above] {$d_3$};
    \end{tikzpicture}
  \end{minipage}
  \hfill
  \begin{minipage}[t]{0.5\linewidth}
    \vspace{1em}
    Lire l'équation réduite de chaque droite.

    \medskip
    \begin{multihideW}[1.5em]
      $d_1$ : $b = 2$, $a = 1$, donc $y = x + 2$.

      \bigskip
      $d_2$ : $b = 4$, $a = -2$, donc $y = -2x + 4$.

      \bigskip
      $d_3$ : $b = -2$. On avance de $3$, on monte de $2$ :

      \smallskip
      $a = \dfrac{2}{3}$, donc $y = \dfrac{2}{3}x - 2$.
    \end{multihideW}
  \end{minipage}
\end{example}

\section{Calcul du coefficient directeur}

\begin{example}
  On reprend la location de vélos : $f(x) = 3x + 5$.

  Pour $2$ heures, on paie $f(2) = \hideM{11}$ €. Pour $5$ heures, on paie $f(5) = \hideM{20}$ €.

  Quand on passe de $2$ à $5$ heures, le prix augmente de \hideM{20 - 11 = 9} € pour \hideM{5 - 2 = 3} heures de plus,
  soit $\dfrac{9}{3} = \hideM{3}$ € par heure supplémentaire. On retrouve le coefficient $a$.
\end{example}

\begin{proposition}{}{}
  \begin{multihide}
    Soit $A(x_A~;~y_A)$ et $B(x_B~;~y_B)$ deux points de la droite d'une fonction affine.

    Le coefficient directeur de cette droite est :
    \[
      a = \frac{y_B - y_A}{x_B - x_A}.
    \]
  \end{multihide}
\end{proposition}

\begin{remark}
  Le coefficient directeur $a$ indique de combien augmente $f(x)$ quand $x$ augmente de $1$.
  Pour la location de vélos, $a = 3$ : le prix augmente de $3$~€ par heure.
\end{remark}

\begin{method}{Déterminer une fonction affine à partir de deux points}{}
  \begin{enumerate}
    \item On calcule $a$ avec la formule $a = \dfrac{y_B - y_A}{x_B - x_A}$.
    \item On calcule $b$ : on remplace $x$ et $y$ par les coordonnées d'un des points dans $y = ax + b$, puis on résout l'équation.
    \item On vérifie avec l'autre point.
  \end{enumerate}
\end{method}

\begin{example}
  Déterminer l'équation réduite de la droite passant par $A(1~;~3)$ et $B(4~;~9)$.

  \begin{multihideW}[1.5em]
    \begin{enumerate}
      \item $a = \dfrac{y_B - y_A}{x_B - x_A} = \dfrac{9 - 3}{4 - 1} = \dfrac{6}{3} = 2$.
      \item Le point $A(1~;~3)$ est sur la droite, donc $3 = 2 \times 1 + b$, soit $3 = 2 + b$, d'où $b = 3 - 2 = 1$.
      \item Vérification avec $B$ : $2 \times 4 + 1 = 9$. \checkmark
    \end{enumerate}
    La droite $(AB)$ a pour équation $y = 2x + 1$.
  \end{multihideW}
\end{example}

\section{Sens de variation}

\begin{proposition}{}{}
  \begin{multihide}
    Soit $f$ une fonction affine définie par $f(x) = ax + b$.
    \begin{itemize}
      \item Si $a > 0$, alors $f$ est \textbf{croissante} sur $\mathbb{R}$ : la droite « monte ».
      \item Si $a < 0$, alors $f$ est \textbf{décroissante} sur $\mathbb{R}$ : la droite « descend ».
      \item Si $a = 0$, alors $f$ est \textbf{constante} sur $\mathbb{R}$ : la droite est horizontale.
    \end{itemize}
  \end{multihide}
\end{proposition}

\begin{center}
  \begin{tabular}{C{0.31\linewidth}C{0.31\linewidth}C{0.31\linewidth}}
    $a > 0$ & $a < 0$ & $a = 0$ \\[0.3em]
    \begin{tikzpicture}[scale=0.45]
      \draw[->] (-2.5,0) -- (2.7,0); \draw[->] (0,-2) -- (0,2.7);
      \hidetk{\draw[blue, line width=1.8pt] (-2.2,-1.6) -- (2.2,2.2);}
    \end{tikzpicture}
    &
    \begin{tikzpicture}[scale=0.45]
      \draw[->] (-2.5,0) -- (2.7,0); \draw[->] (0,-2) -- (0,2.7);
      \hidetk{\draw[blue, line width=1.8pt] (-2.2,2.2) -- (2.2,-1.6);}
    \end{tikzpicture}
    &
    \begin{tikzpicture}[scale=0.45]
      \draw[->] (-2.5,0) -- (2.7,0); \draw[->] (0,-2) -- (0,2.7);
      \hidetk{\draw[blue, line width=1.8pt] (-2.2,1.2) -- (2.2,1.2);}
    \end{tikzpicture}
    \\
    \begin{tikzpicture}[scale=0.8]
      \tkzTabInit[lgt=1,espcl=2.2]{$x$/1,$f$/1.5}{$-\infty$,$+\infty$}
      \hidetk{\tkzTabVar{-/ ,+/ }}
    \end{tikzpicture}
    &
    \begin{tikzpicture}[scale=0.8]
      \tkzTabInit[lgt=1,espcl=2.2]{$x$/1,$f$/1.5}{$-\infty$,$+\infty$}
      \hidetk{\tkzTabVar{+/ ,-/ }}
    \end{tikzpicture}
    &
    \begin{tikzpicture}[scale=0.8]
      \tkzTabInit[lgt=1,espcl=2.2]{$x$/1,$f$/1.5}{$-\infty$,$+\infty$}
      \hidetk{\tkzTabVar{+/ $b$,+/ $b$}}
    \end{tikzpicture}
  \end{tabular}
\end{center}

\begin{examples}
  Donner le sens de variation de chaque fonction.

  \medskip
  \begin{minipage}[t]{0.31\linewidth}
    $f(x) = 3x - 7$
    \begin{multihideW}[1.5em]
      $a = 3 > 0$, donc $f$ est croissante sur $\mathbb{R}$.
    \end{multihideW}
  \end{minipage}
  \hfill
  \begin{minipage}[t]{0.31\linewidth}
    $g(x) = -2x + 5$
    \begin{multihideW}[1.5em]
      $a = -2 < 0$, donc $g$ est décroissante sur $\mathbb{R}$.
    \end{multihideW}
  \end{minipage}
  \hfill
  \begin{minipage}[t]{0.31\linewidth}
    $h(x) = 4 - x$
    \begin{multihideW}[1.5em]
      $a = -1 < 0$, donc $h$ est décroissante sur $\mathbb{R}$.
    \end{multihideW}
  \end{minipage}
\end{examples}

\newpage

\section{Signe d'une fonction affine}

\begin{minipage}[c]{0.56\linewidth}
  \begin{example*}
    Soit $f(x) = 2x - 6$.

      \begin{itemize}
        \item On cherche quand $f(x) = 0$ :

        $2x - 6 = 0 \iff \hideM{2x = 6} \iff \hideM{x = 3}$.
        \item $a = 2 > 0$, donc $f$ est \hide{croissante} : $f(x)$ est \hide{négatif} avant $3$ et \hide{positif} après $3$.
      \end{itemize}

      \medskip
      Tableau de signes de $f$ :

      \begin{center}
        \begin{tikzpicture}
          \tkzTabInit[lgt=1.5,espcl=2]{$x$/1,$f(x)$/1}{$-\infty$,$\hideM{3}$,$+\infty$}
          \hidetk{\tkzTabLine{,-,z,+,}}
        \end{tikzpicture}
      \end{center}
  \end{example*}
\end{minipage}
\hfill
\begin{minipage}[c]{0.4\linewidth}
  \centering
  \begin{tikzpicture}[scale=0.8]
    \repere{-1}{5}{-3}{3}
    \draw[very thick, red, domain=1.5:3] plot (\x,{2*\x-6});
    \draw[very thick, green!50!black, domain=3:4.5] plot (\x,{2*\x-6});
    \fill (3,0) circle (3pt);
  \end{tikzpicture}
\end{minipage}

\begin{proposition}{}{}
  \begin{multihide}
    Soit $f(x) = ax + b$ avec $a \neq 0$. La fonction $f$ s'annule en $x = -\dfrac{b}{a}$, et :% (solution de $ax + b = 0$), et :

    \medskip
    \begin{minipage}{0.48\linewidth}
      \centering
      Si $a > 0$ :
      \vspace{0.5em}

      \begin{tikzpicture}[scale=0.85]
        \tkzTabInit[lgt=2,espcl=1.8]{$x$/1,$ax+b$/1}{$-\infty$,$-\frac{b}{a}$,$+\infty$}
        \tkzTabLine{,-,z,+,}
      \end{tikzpicture}
    \end{minipage}
    \hfill
    \begin{minipage}{0.48\linewidth}
      \centering
      Si $a < 0$ :
      \vspace{0.5em}

      \begin{tikzpicture}[scale=0.85]
        \tkzTabInit[lgt=2,espcl=1.8]{$x$/1,$ax+b$/1}{$-\infty$,$-\frac{b}{a}$,$+\infty$}
        \tkzTabLine{,+,z,-,}
      \end{tikzpicture}
    \end{minipage}

    %\medskip
    %À retenir : à droite du $0$, $ax + b$ a le \textbf{signe de $a$}.
  \end{multihide}
\end{proposition}

\begin{method}{Dresser le tableau de signes de $ax + b$}{}
  \begin{enumerate}
    \item On résout l'équation $ax + b = 0$.
    \item On regarde le signe de $a$.
    \item On remplit le tableau : $0$ sous la solution, signe de $a$ à droite, signe contraire à gauche.
  \end{enumerate}
\end{method}

\begin{example*}
  Dresser le tableau de signes de $g(x) = -3x + 12$.
  \begin{multihideW}
    \begin{minipage}{0.5\linewidth}
      $-3x + 12 = 0 \iff -3x = -12 \iff x = 4$.

      \smallskip
      $a = -3 < 0$, donc $a$ est négatif.
    \end{minipage}
    \hfill
    \begin{minipage}{0.45\linewidth}
      \begin{tikzpicture}
        \tkzTabInit[lgt=1.5,espcl=1.8]{$x$/1,$g(x)$/1}{$-\infty$,$4$,$+\infty$}
        \tkzTabLine{,+,z,-,}
      \end{tikzpicture}
    \end{minipage}
  \end{multihideW}
\end{example*}

\vfill

\begin{savoir}{}{}
\begin{itemize}
  \item Reconnaître une fonction affine et identifier $a$ et $b$,
  \item Modéliser une situation par une fonction affine,
  \item Tracer une droite à partir de son équation réduite,
  \item Lire graphiquement l'équation réduite d'une droite,
  \item Calculer le coefficient directeur à partir de deux points,
  \item Donner le sens de variation selon le signe de $a$,
  \item Dresser le tableau de signes de $ax + b$.
\end{itemize}
\end{savoir}

\end{document}
